\section{Обговорення}
\label{sec:discussion}

Основна перевага запропонованого
двомасштабного підходу полягає
в можливості поєднувати детальне
стохастичне моделювання всередині
спільнот з агрегованим представленням
зв'язків між ними. Це відповідає
задачам, у яких дані про контактні
мережі доступні з різним рівнем
деталізації. Наприклад, внутрішню
структуру спільнот можна відновлювати
або генерувати за відомими
характеристиками, тоді як
міжспільнотні зв'язки можуть
бути представлені лише загальними
потоками переміщень, отриманими
з транспортних чи міграційних даних.
На відміну від повномережевого
моделювання, \TS{} не потребує
явного визначення вузлів,
між якими відбуваються
міжспільнотні контакти.
Водночас, на відміну від
повністю агрегованих моделей,
він зберігає індивідуальну
динаміку зараження та вплив
внутрішньої топології.

Експерименти продемонстрували,
що таке спрощення може зберігати
високу точність часових характеристик
поширення інфекції. У спільнотах
зі швидким внутрішнім поширенням
розподіли часу прибуття у \Full{}
і \TS{} виявилися близькими,
причому це узгодження зберігалося
навіть після дев'яти послідовних
міжспільнотних переходів.
Отже, відсутність інформації
про конкретні вузли-мости
не обов'язково призводить
до суттєвої похибки
в загальній динаміці поширення.
Це робить агреговані інтерфейси
перспективним засобом
моделювання мереж спільнот,
коли основний інтерес становлять
часи досягнення спільнот
та їхні розподіли.

Точність такого наближення
можна пояснити співвідношенням
часових масштабів.
Якщо інфекція швидко охоплює
вузли спільноти порівняно
з очікуванням міжспільнотного
передавання, конкретне
розташування мостів має
порівняно невеликий вплив
на час експорту. Натомість
за повільнішого та нерівномірнішого
внутрішнього поширення
стан окремих вузлів-мостів
може залишатися важливим
протягом значної частини
часу очікування переходу.
Саме тому структурні
характеристики, зокрема ASPL,
не повністю пояснюють
спостережені розбіжності.
Сильніший зв'язок точності
з показником
$\rho=\beta_{\mathrm{out}}\sigma_T$
підтримує пояснення через
часову нерівномірність
внутрішнього зараження.
Це відкриває можливість
оцінювати доцільність
агрегації не лише за
топологією мережі,
а й за динамікою процесу.

Крім того, можна продовжувати покращувати деталізацію включивши локальні параметри спільнот до моделі (як от ASPL, власні значення чи інші), продовження зниження різниці між методами, що є ідеєю для подальших досліджень.

Для моделювання поширення
через багато спільнот
важливим є також характер
накопичення розбіжностей.
Експерименти показали,
що абсолютна різниця
між часами прибуття може
збільшуватися вздовж ланцюга,
тоді як відносний зсув
часто стабілізується
або змінюється повільно.
Отже, похибка агрегованого
моделювання не обов'язково
зростає пропорційно
кількості міжспільнотних
переходів. Водночас
сильніше зміщення
сукупної тривалості
наступних переходів
порівняно з першим
вказує на можливі
додаткові ефекти
після імпорту заражень.
Одним із пояснень
є виникнення надлишкових
імпортів, які можуть
створювати додаткові
осередки інфекції
та прискорювати
подальше поширення.
Окремий внесок цього
механізму потребує
додаткової перевірки.

Теоретичний аналіз
дозволяє чітко визначити,
яке наближення реалізує
двомасштабний підхід.
Ідеальний процес
еквівалентний точному
моделюванню на сурогатній
мережі з рівномірно
розподіленими
міжспільнотними зв'язками.
Це дає математичне
обґрунтування агрегованого
представлення та пояснює,
чому його результати
можуть відрізнятися
від повномережевого
моделювання.
Зокрема, близькість
часових розподілів
не гарантує відтворення
конкретних маршрутів
зараження. Тому вибір
рівня деталізації
має визначатися
метою дослідження:
для оцінювання часових
характеристик агреговане
представлення може бути
достатнім, тоді як
аналіз фізичних шляхів
передавання потребує
інформації про відповідні
контакти.

Додаткові можливості
методу пов'язані
з його обчислювальною
структурою. Незалежні
локальні симулятори
можуть виконуватися
паралельно між моментами
синхронізації, а окремі
спільноти можуть мати
різну внутрішню структуру
та параметри моделювання.
Крім того, заміна
явних міжспільнотних
ребер агрегованими вагами
може зменшувати обсяг
даних про контакти.
Проте фактичний виграш
у швидкодії та пам'яті
залежить від структури
мережі й витрат на
синхронізацію.

% Подальший розвиток
% підходу пов'язаний
% насамперед із визначенням
% умов, за яких можна
% використовувати агрегацію
% без суттєвої втрати
% точності. Перспективним
% є адаптивне представлення
% міжспільнотних інтерфейсів:
% застосування агрегованих
% зв'язків там, де їх
% достатньо для відтворення
% часової динаміки,
% та збереження детальних
% контактів у спільнотах,
% де стан вузлів-мостів
% має значний вплив.
% Окремої перевірки
% потребує застосування
% методу до мереж
% із наближено відтвореною
% внутрішньою структурою,
% оскільки в наведених
% експериментах внутрішні
% графи обох моделей
% були однаковими.
% Також доцільно
% перевірити можливість
% оцінювання співвідношення
% часових масштабів
% без попереднього
% повномережевого
% моделювання та дослідити
% інші епідемічні процеси.









\section{Висновки}
\label{sec:conclusions}

У роботі досліджено
двомасштабний метод
стохастичного моделювання
поширення інфекції
в мережах спільнот.
Метод поєднує детальне
моделювання внутрішніх
контактних мереж
за алгоритмом Gillespie
з агрегованими
міжспільнотними
інтенсивностями.
Це дозволяє використовувати
локальні контактні мережі
разом із загальними
даними про зв'язки
між спільнотами,
без необхідності
явного представлення
кожного міжспільнотного
контакту. Модульна
структура також
створює можливість
паралельного виконання
локальних симуляторів.

Теоретично доведено
еквівалентність ідеального
двомасштабного процесу
точному моделюванню
на сурогатній мережі
з рівномірними
міжспільнотними
зв'язками. Проаналізовано
механізми розбіжностей
між агрегованим
і повномережевим
моделюванням,
зокрема наслідки
втрати інформації
про конкретні
вузли-мости. Крім того, доведено, що за
рівномірного випадкового розміщення
вузлів-мостів агрегація спричиняє
стохастично не пізніший перший
експорт інфекції в ідеальній моделі.

Експериментально
підтверджено здатність
\TS{} відтворювати
близькі до \Full{}
часові розподіли
поширення інфекції
за швидкого внутрішнього
зараження спільнот.
Це узгодження
зберігалося навіть
після дев'яти
послідовних
міжспільнотних
переходів.
Встановлено, що
абсолютна розбіжність
часів прибуття
може накопичуватися
вздовж ланцюга,
тоді як відносна
часто стабілізується
або змінюється
повільно.

Досліджено вплив
внутрішньої структури,
розміру спільнот
і кількості
міжспільнотних
ребер на точність
агрегованого
моделювання.
Показано, що
співвідношення
часових масштабів
$\rho=\beta_{\mathrm{out}}\sigma_T$
має сильніший
спостережений зв'язок
із розбіжностями,
ніж ASPL, і є
перспективною
характеристикою
для визначення
умов застосовності
методу.

У наших тестах найменші розбіжності спостерігалися
за швидкого внутрішнього зараження. У всіх 18
конфігураціях із $\rho<0.25$ нормована відстань
Вассерштайна не перевищувала 4.5\%.
Водночас близькі часові розподіли не гарантували
збереження допустимих маршрутів передавання.






\clearpage

